% ─────────────────────────────────────────────────────────────────────────────
\section{Le groupement par mat\'eriau~: \'eviter les calculs en double}
% ─────────────────────────────────────────────────────────────────────────────

\begin{maitre}
Si ton stock contient 50~poutres C24 de 63\texttimes{}175, dois-tu
les calculer 50 fois~?
\end{maitre}

\begin{apprenti}
Non, elles ont toutes les m\^emes propri\'et\'es m\'ecaniques~!
Il suffit de calculer une fois et de copier le r\'esultat aux 50~produits.
\end{apprenti}

\begin{maitre}
C'est exactement ce que fait le code.
\end{maitre}

\begin{lstlisting}[style=python, caption={Groupement par \py{id\_config\_materiau} --- extrait de \fichier{moteur.py}}]
# groupes : clé = id matériau, valeur = liste de produits identiques
groupes: dict[str, list[ProduitValide]] = {}

for produit in produits_valides:
    # setdefault : si la clé n'existe pas, crée une liste vide d'abord
    groupes.setdefault(produit.id_config_materiau, []).append(produit)

# Résultat :
# groupes = {
#   "C24-063-175-6000": [prod_1, prod_2, prod_3],  # 3 refs 63x175
#   "C24-075-225-7000": [prod_4],                   # 1 ref  75x225
# }

for id_mat, produits_groupe in groupes.items():
    p_rep = produits_groupe[0]     # représentant du groupe pour le calcul

    # ... calculs EC5 (une seule fois pour tout le groupe) ...

    # Réplication : copier le résultat à chaque produit du groupe
    for produit in produits_groupe:
        # dataclasses.replace() crée une copie avec un champ modifié
        r_produit = replace(r_base, id_produit=produit.id_produit)
        tous.append(r_produit)
\end{lstlisting}